% ============================================================
\chapter{Mesures Produits et Th\'eor\`eme de Fubini-Tonelli}
\label{ch:fubini}
% ============================================================

\section{Introduction}

Le calcul d'int\'egrales multiples est l'une des op\'erations les plus fr\'equentes
en analyse et en probabilit\'es. Les th\'eor\`emes de Fubini et de Tonelli fournissent
les conditions sous lesquelles une int\'egrale double peut \^etre calcul\'ee comme
une int\'egrale it\'er\'ee\index{int\'egrale it\'er\'ee}, et dans n'importe quel ordre.

\begin{intuition}[title=\textbf{Pourquoi Fubini-Tonelli~?}]
Calculer $\iint f(x,y)\, d(x,y)$ directement sur un produit est souvent difficile.
Fubini et Tonelli permettent de ramener ce calcul \`a des int\'egrales
successives~: d'abord int\'egrer en $y$ pour chaque $x$ fix\'e (ou vice versa).
Tonelli traite le cas des fonctions positives (sans hypoth\`ese d'int\'egrabilit\'e),
tandis que Fubini s'applique aux fonctions int\'egrables.
\end{intuition}

% ------------------------------------------------------------
\section{$\sigma$-alg\`ebre produit}
% ------------------------------------------------------------

\begin{definition}[$\sigma$-alg\`ebre produit]
Soient $(X, \mathcal{A})$ et $(Y, \mathcal{B})$ deux espaces mesurables. La
\emph{$\sigma$-alg\`ebre produit}\index{$\sigma$-alg\`ebre produit}
$\mathcal{A} \otimes \mathcal{B}$ est la $\sigma$-alg\`ebre sur $X \times Y$
engendr\'ee par les \emph{rectangles mesurables}~:
\[
    \mathcal{A} \otimes \mathcal{B}
    = \sigma\bigl(\{A \times B : A \in \mathcal{A},\, B \in \mathcal{B}\}\bigr).
\]
\end{definition}

\begin{definition}[Sections]
Pour $E \subset X \times Y$ et $x \in X$, $y \in Y$, on d\'efinit les
\emph{sections}\index{section}~:
\[
    E_x = \{y \in Y : (x,y) \in E\}, \qquad
    E^y = \{x \in X : (x,y) \in E\}.
\]
Pour une fonction $f : X \times Y \to \overline{\R}$, on d\'efinit les sections~:
$f_x(y) = f(x,y)$ et $f^y(x) = f(x,y)$.
\end{definition}

\begin{proposition}[Mesurabilit\'e des sections]
\label{prop:sections_mesurables}
Soit $E \in \mathcal{A} \otimes \mathcal{B}$. Alors pour tout $x \in X$, la section
$E_x \in \mathcal{B}$, et pour tout $y \in Y$, la section $E^y \in \mathcal{A}$.

Si $f : (X \times Y, \mathcal{A} \otimes \mathcal{B}) \to (\overline{\R}, \mathcal{B}(\overline{\R}))$
est mesurable, alors $f_x$ est $\mathcal{B}$-mesurable et $f^y$ est
$\mathcal{A}$-mesurable.
\end{proposition}

\begin{proof}
La collection $\mathcal{C} = \{E \in \mathcal{A} \otimes \mathcal{B} : E_x \in \mathcal{B}
\text{ pour tout } x\}$ est une $\sigma$-alg\`ebre contenant les rectangles mesurables
(car $(A \times B)_x = B$ si $x \in A$ et $\varnothing$ sinon), donc
$\mathcal{C} = \mathcal{A} \otimes \mathcal{B}$.
\end{proof}

% ------------------------------------------------------------
\section{Mesure produit}
% ------------------------------------------------------------

\begin{theoreme}[Existence et unicit\'e de la mesure produit]
\label{thm:mesure_produit}
\index{mesure produit}
Soient $(X, \mathcal{A}, \mu)$ et $(Y, \mathcal{B}, \nu)$ deux espaces mesur\'es
$\sigma$-finis. Il existe une unique mesure $\mu \otimes \nu$ sur
$(X \times Y, \mathcal{A} \otimes \mathcal{B})$ telle que~:
\[
    (\mu \otimes \nu)(A \times B) = \mu(A) \cdot \nu(B),
    \qquad \forall A \in \mathcal{A},\, \forall B \in \mathcal{B}.
\]
De plus, pour tout $E \in \mathcal{A} \otimes \mathcal{B}$~:
\[
    (\mu \otimes \nu)(E) = \int_X \nu(E_x)\, d\mu(x)
    = \int_Y \mu(E^y)\, d\nu(y).
\]
\end{theoreme}

\begin{proof}[Id\'ee de la preuve]
On d\'efinit $\rho(E) = \int_X \nu(E_x)\, d\mu(x)$. La mesurabilit\'e de
$x \mapsto \nu(E_x)$ se v\'erifie par un argument de classe monotone. La
$\sigma$-additivit\'e de $\rho$ r\'esulte du th\'eor\`eme de convergence monotone.
Sur les rectangles, $\rho(A \times B) = \int_X \nu(B) \ind_A(x)\, d\mu(x) =
\mu(A) \nu(B)$. L'unicit\'e d\'ecoule de l'unicit\'e des mesures sur un $\pi$-syst\`eme
(les rectangles forment un $\pi$-syst\`eme) et de la $\sigma$-finitude.
\end{proof}

\begin{attention}[title=\textbf{Hypoth\`ese de $\sigma$-finitude}]
La $\sigma$-finitude est essentielle. Sans elle, la mesure produit peut ne pas \^etre
unique, et les th\'eor\`emes de Fubini et Tonelli peuvent \^etre en d\'efaut.
\end{attention}

% ------------------------------------------------------------
\section{Th\'eor\`eme de Tonelli}
% ------------------------------------------------------------

\begin{theoreme}[Tonelli]
\label{thm:tonelli}
\index{th\'eor\`eme de Tonelli}
Soient $(X, \mathcal{A}, \mu)$ et $(Y, \mathcal{B}, \nu)$ deux espaces mesur\'es
$\sigma$-finis et soit $f : X \times Y \to [0, +\infty]$ une fonction
$(\mathcal{A} \otimes \mathcal{B})$-mesurable. Alors~:
\begin{enumerate}
    \item[(i)] pour $\mu$-presque tout $x$, la fonction $y \mapsto f(x,y)$ est
    $\mathcal{B}$-mesurable et $x \mapsto \int_Y f(x,y)\, d\nu(y)$ est
    $\mathcal{A}$-mesurable\,;
    \item[(ii)] de m\^eme en intervertissant les r\^oles de $x$ et $y$\,;
    \item[(iii)] on a l'\'egalit\'e~:
    \[
        \int_{X \times Y} f\, d(\mu \otimes \nu)
        = \int_X \left(\int_Y f(x,y)\, d\nu(y)\right) d\mu(x)
        = \int_Y \left(\int_X f(x,y)\, d\mu(x)\right) d\nu(y).
    \]
\end{enumerate}
\end{theoreme}

\begin{proof}
Pour $f = \ind_E$ avec $E \in \mathcal{A} \otimes \mathcal{B}$, le r\'esultat est le
th\'eor\`eme~\ref{thm:mesure_produit}. Par lin\'earit\'e, il s'\'etend aux fonctions
\'etag\'ees positives. Le cas g\'en\'eral s'obtient par le th\'eor\`eme de convergence
monotone appliqu\'e \`a une suite croissante de fonctions \'etag\'ees convergeant vers $f$.
\end{proof}

% ------------------------------------------------------------
\section{Th\'eor\`eme de Fubini}
% ------------------------------------------------------------

\begin{theoreme}[Fubini]
\label{thm:fubini}
\index{th\'eor\`eme de Fubini}
Soient $(X, \mathcal{A}, \mu)$ et $(Y, \mathcal{B}, \nu)$ deux espaces mesur\'es
$\sigma$-finis et soit $f \in L^1(X \times Y, \mu \otimes \nu)$. Alors~:
\begin{enumerate}
    \item[(i)] pour $\mu$-presque tout $x$, $f_x \in L^1(Y, \nu)$\,;
    \item[(ii)] la fonction $x \mapsto \int_Y f(x,y)\, d\nu(y)$ (d\'efinie p.p.)
    est dans $L^1(X, \mu)$\,;
    \item[(iii)] on a~:
    \[
        \int_{X \times Y} f\, d(\mu \otimes \nu)
        = \int_X \left(\int_Y f(x,y)\, d\nu(y)\right) d\mu(x)
        = \int_Y \left(\int_X f(x,y)\, d\mu(x)\right) d\nu(y).
    \]
\end{enumerate}
\end{theoreme}

\begin{proof}
On \'ecrit $f = f^+ - f^-$. Par Tonelli, $\int f^+\, d(\mu \otimes \nu)$ et
$\int f^-\, d(\mu \otimes \nu)$ sont finis (car $f \in L^1$). Le th\'eor\`eme de Tonelli
appliqu\'e \`a $f^+$ et $f^-$ s\'epar\'ement donne les int\'egrales it\'er\'ees pour chacun.
En particulier, $\int_Y f^+(x,y)\, d\nu(y) < \infty$ et
$\int_Y f^-(x,y)\, d\nu(y) < \infty$ pour $\mu$-presque tout $x$, donc
$f_x \in L^1(\nu)$ p.p. On conclut par diff\'erence.
\end{proof}

\begin{formules}[title=\textbf{Fubini-Tonelli~: mode d'emploi}]
\begin{enumerate}
    \item \textbf{Tonelli}~: $f \geq 0$ mesurable $\implies$ on peut intervertir
    $\int\int$ librement (m\^eme si $= +\infty$).
    \item \textbf{Fubini}~: $f \in L^1(\mu \otimes \nu)$ $\implies$ on peut intervertir.
    \item \textbf{Strat\'egie}~: pour v\'erifier $f \in L^1$, appliquer d'abord Tonelli
    \`a $\abs{f}$.
\end{enumerate}
\end{formules}

% ------------------------------------------------------------
\section{Applications}
% ------------------------------------------------------------

\subsection{Convolution}

\begin{definition}[Convolution]
\index{convolution}
Soient $f, g \in L^1(\R^n, \lambda)$ o\`u $\lambda$ est la mesure de Lebesgue.
La \emph{convolution} de $f$ et $g$ est~:
\[
    (f * g)(x) = \int_{\R^n} f(x - y)\, g(y)\, dy.
\]
\end{definition}

\begin{proposition}
Si $f, g \in L^1(\R^n)$, alors $f * g$ est d\'efinie p.p., $f * g \in L^1(\R^n)$, et
$\norm{f * g}_{L^1} \leq \norm{f}_{L^1} \norm{g}_{L^1}$.
\end{proposition}

\begin{proof}
Par Tonelli, $\int_{\R^n}\!\int_{\R^n} \abs{f(x-y)} \abs{g(y)}\, dy\, dx
= \norm{f}_{L^1} \norm{g}_{L^1} < \infty$ (par translation). Par Fubini,
l'int\'egrale int\'erieure est finie pour presque tout $x$, et le r\'esultat suit.
\end{proof}

\subsection{Formule de changement de variable}

\begin{theoreme}[Changement de variable]
\index{changement de variable}
Soit $\varphi : U \to V$ un $C^1$-diff\'eomorphisme entre deux ouverts de $\R^n$.
Pour toute fonction mesurable $f : V \to [0, +\infty]$ (ou $f \in L^1(V)$)~:
\[
    \int_V f(y)\, dy = \int_U f(\varphi(x))\, \abs{\det D\varphi(x)}\, dx.
\]
\end{theoreme}

\begin{exemple}[Coordonn\'ees polaires]
Pour $\varphi(r, \theta) = (r\cos\theta, r\sin\theta)$, on a
$\abs{\det D\varphi} = r$, d'o\`u~:
\[
    \iint_{\R^2} f(x,y)\, dx\, dy = \int_0^{\infty} \int_0^{2\pi} f(r\cos\theta, r\sin\theta)\, r\, d\theta\, dr.
\]
\end{exemple}

\begin{exemple}[Int\'egrale de Gauss]
En utilisant Tonelli et les coordonn\'ees polaires~:
\[
    \left(\int_{-\infty}^{\infty} e^{-x^2}\, dx\right)^2
    = \int_{\R^2} e^{-(x^2+y^2)}\, dx\, dy
    = \int_0^{\infty} 2\pi r\, e^{-r^2}\, dr = \pi.
\]
Donc $\int_{-\infty}^{\infty} e^{-x^2}\, dx = \sqrt{\pi}$.
\end{exemple}

% ------------------------------------------------------------
\section{Compl\'ements~: produits infinis}
% ------------------------------------------------------------

\begin{theoreme}[Produit d\'enombrable de mesures de probabilit\'e]
Soit $(X_n, \mathcal{A}_n, \mu_n)_{n \geq 1}$ une suite d'espaces de probabilit\'e.
Il existe une unique probabilit\'e $\PP$ sur
$\left(\prod_{n \geq 1} X_n,\, \bigotimes_{n \geq 1} \mathcal{A}_n\right)$ telle que~:
\[
    \PP\!\left(\prod_{n=1}^{N} A_n \times \prod_{n > N} X_n\right)
    = \prod_{n=1}^{N} \mu_n(A_n)
\]
pour tout $N$ et tout choix de $A_n \in \mathcal{A}_n$. C'est le
\emph{th\'eor\`eme d'Ionescu-Tulcea} ou le th\'eor\`eme d'extension de Kolmogorov
dans le cas produit.
\end{theoreme}

% ------------------------------------------------------------
\section{Exercices}
% ------------------------------------------------------------

\begin{exercice}
Calculer $\int_0^1 \int_0^1 \frac{x^2 - y^2}{(x^2 + y^2)^2}\, dx\, dy$ et
$\int_0^1 \int_0^1 \frac{x^2 - y^2}{(x^2 + y^2)^2}\, dy\, dx$. Commenter.
\emph{Indication}~: $f \notin L^1$.
\end{exercice}

\begin{exercice}
Soit $f : [0,1] \times [0,1] \to \R$ d\'efinie par $f(x,y) = \frac{xy}{(x^2+y^2)^2}$.
V\'erifier que $f \notin L^1([0,1]^2)$ et que les int\'egrales it\'er\'ees donnent
des r\'esultats diff\'erents.
\end{exercice}

\begin{exercice}
En utilisant Fubini, montrer que pour $f \in L^1([0,\infty), \lambda)$~:
\[
    \int_0^{\infty} \frac{1}{x} \int_0^x f(t)\, dt\, dx
    = \int_0^{\infty} f(t) \int_t^{\infty} \frac{1}{x}\, dx\, dt
\]
et expliquer pourquoi cette \'egalit\'e est formelle (les deux c\^ot\'es peuvent diverger).
\end{exercice}

\begin{exercice}
Montrer que $\int_0^{\infty} \frac{e^{-ax} - e^{-bx}}{x}\, dx = \ln(b/a)$ pour
$0 < a < b$, en utilisant Fubini.
\emph{Indication}~: \'ecrire $\frac{e^{-ax} - e^{-bx}}{x} = \int_a^b e^{-tx}\, dt$.
\end{exercice}

\begin{exercice}[Convolution et loi de probabilit\'e]
Soient $X$ et $Y$ deux variables al\'eatoires r\'eelles ind\'ependantes de densit\'es
respectives $f_X$ et $f_Y$. Montrer que $Z = X + Y$ admet la densit\'e
$f_Z = f_X * f_Y$. Appliquer au cas $X, Y \sim \mathcal{N}(0,1)$.
\end{exercice}

\begin{exercice}
En utilisant les coordonn\'ees polaires en dimension $n$, calculer le volume de la
boule unit\'e $B_n = \{x \in \R^n : \norm{x} \leq 1\}$ et montrer que~:
\[
    \lambda_n(B_n) = \frac{\pi^{n/2}}{\Gamma(n/2 + 1)}.
\]
\end{exercice}
